\chapter{Modul Master Data}

Modul Master Data merupakan fondasi sistem. Seluruh transaksi penjualan, pembelian, dan pergudangan bergantung pada integritas data yang dimasukkan ke dalam modul ini.

\section{Manajemen Produk (\textit{Product})}

Modul \textit{Product} menyimpan daftar inventaris, katalog produk untuk \textit{website} publik, dan pengaturan harga.

\subsection{Alur Kerja Pembuatan Produk Baru}
\begin{enumerate}
    \item Masuk ke menu navigasi \textbf{Master Data} $>$ \textbf{Products}.
    \item Klik tombol \textbf{New Product} di sudut kanan atas halaman.
    \item Isi formulir informasi dasar produk:
    \begin{itemize}
        \item \textbf{Nama Produk}: Tulis nama jelas (contoh: "Sepeda Lipat BMX").
        \item \textbf{SKU}: Isi dengan kode unik barang.
        \item \textbf{Harga Beli}: Modal barang (hanya bisa dilihat oleh \textit{Owner}/\textit{Super Admin}).
        \item \textbf{Margin Penjualan (\%)}: Sistem otomatis akan menghitung \textbf{Harga Jual} berdasarkan persentase margin ini.
    \end{itemize}
    \item Unggah foto produk pada \textit{field} \textbf{Foto Utama} dan \textbf{Galeri Foto}.
    \item Atur visibilitas produk dengan menghidupkan \textit{toggle} \textbf{Tampilkan di Website}.
    \item Klik tombol \textbf{Create} untuk menyimpan data.
\end{enumerate}

\begin{figure}[H]
    \centering
    \includegraphics[width=0.9\textwidth]{placeholder_form_produk.png}
    \caption{Formulir Penambahan Produk Baru}
    \label{figure:form_produk}
\end{figure}

\section{Manajemen Customer \& Supplier}

Sistem membedakan profil rekanan bisnis melalui menu \textbf{Customers} dan \textbf{Suppliers}.

\subsection{Pengelolaan Data Customer}
\begin{enumerate}
    \item Klik menu \textbf{Customers}.
    \item Klik \textbf{New Customer}.
    \item Isi data wajib meliputi \textbf{Nama Lengkap}, \textbf{Nomor WhatsApp}, \textbf{Email}, dan \textbf{Alamat Pengiriman}.
    \item Status \textbf{Is Active} harus disetel hidup agar nama pelanggan muncul di menu pembuatan \textit{Quotation}.
    \item Klik \textbf{Create}.
\end{enumerate}

\section{Warehouse (Gudang)}

Gudang menyimpan informasi lokasi fisik tempat penyimpanan stok. Manajemen yang tidak memiliki \textit{Role} "Gudang" dapat melihat tetapi terbatas dalam mengeksekusi mutasi.

\subsection{Skenario Dunia Nyata: Penambahan Cabang Gudang}
Perusahaan membuka cabang baru di Surabaya. \textit{Super Admin} membuat profil Gudang:
\begin{enumerate}
    \item Buka menu \textbf{Warehouses}.
    \item Klik \textbf{New Warehouse}.
    \item Isi \textbf{Nama Gudang} dengan "Gudang Cabang Surabaya" dan \textbf{Alamat Lengkap}.
    \item Klik \textbf{Create}. Setelah itu, opsi Gudang Surabaya akan muncul saat staf membuat dokumen \textit{Delivery Order} (Surat Jalan).
\end{enumerate}
